\chapter*{前言}

第一册讲清楚了一个 \cpp 程序怎样变成进程、指令、寄存器、栈帧、二进制接口和可测量的性能现象。第二册不能直接跳到 AI 算子，因为算子开发不是只会写矩阵乘就够了。真正能在本地 Linux 上跑起来的 CPU 推理引擎，先要理解现代多核心 CPU、缓存层级、TLB、预取、NUMA、线程、锁、原子操作、内存模型、性能计数器和软件优化边界。没有这些底座，后面的矩阵乘、量化、Attention 和推理图都只会变成孤立技巧。

本册的路径因此重新定位为三步。第一步是现代 CPU 性能工程：在 Linux 上认识多核机器，知道数据怎样穿过 cache line、TLB、内存带宽和执行单元，知道 benchmark 怎样写才可信。第二步是并发和软件优化：理解线程池、任务切分、锁、原子、false sharing、内存序和缓存友好布局，能解释为什么加线程有时更慢。第三步才进入 AI：从张量、模型、推理图、算子、参数、激活和量化格式讲起，让没有 AI 背景的读者也能逐步读懂模型推理，然后开始实现 CPU 算子和本地推理引擎。

本册有一个贯穿全书的大作业：实现一个 Linux 本地 CPU 版本的量化推理引擎。这个项目不是最后一章临时出现，而是从第一章开始逐步搭建。前面章节先建立工程边界和测量框架，中间章节补线程池、张量结构、内存布局和基础算子，后面章节加入量化权重、模型加载、推理图执行、端到端 benchmark 和正确性验证。最终目标是在本地 Linux 机器上，不依赖云服务，加载一个小型量化模型并完成推理。

\begin{keyidea}
第二册的核心不是“AI 很快”，而是“怎样在 Linux 上把现代 CPU、多核并发、内存层级和 AI 算子组织成一个能运行、能验证、能测量、能继续优化的量化推理系统”。
\end{keyidea}

本册仍然保持第一册的写法：原理先行，代码作为证据；实验必须说明输入、环境、命令和误差；源码阅读抓数据结构、热路径和性能边界。读者读完后，不应该只会背 SIMD、GEMM、Attention 这些名词，而应该能从零解释一个本地 CPU 推理引擎需要哪些模块，能写出可测试的 C++ 组件，能用 Linux 工具定位瓶颈，并能把量化模型跑通。
